\subsection{Applying the Moral Ecosystem Approach to AI Development}
\label{sec:5.4.1}
The reason to borrow the structure is that it locates influences a design review does not look for. Each ring supplies a different kind of pressure on what a system ends up valuing, and none of them is visible in a training run. The team and its data are the innermost, and a team's composition decides what the system is built to notice. Around that sit the connections between the parts of an organization, since a system whose principles are read one way by the engineers and another by the policy staff has no consistent principles. Further out are the structures the work never touches and answers to anyway — regulation, industry standards, the incentives a firm's revenue actually runs on. Outside that is the cultural and economic context, which decides what counts as an obvious product, and around all of it is time, which keeps moving after release.

Two things follow. No ring is sufficient: a diverse team still answers to its employer's incentives, and a well-drafted regulation still depends on somebody interpreting it consistently. And the direction of the whole is set by whichever ring is pushing hardest, which is rarely the one a development team is looking at.
